\documentclass[11pt]{article}

\usepackage[margin=1in]{geometry}
\usepackage{amsmath,amssymb}
\usepackage{microtype}
\usepackage{hyperref}
\usepackage{xcolor}
\usepackage{enumitem}
\usepackage{array}

\hypersetup{
  colorlinks=true,
  linkcolor=black,
  urlcolor=blue!60!black,
}

\title{AniSync App: Algorithm and Recommendation Flow}
\author{Tony Dong \and Edward Kim \and Jakob Rees \and Vincent Wang \and Jack Xiong}
\date{April 2026}

\begin{document}
\maketitle

\section{Setup and notation}

A room contains $n \geq 2$ users. We represent each user $u \in \{1, \dots, n\}$
as a pair
\[
  P_u \;=\; (L_u,\, T_u),
\]
where:
\begin{itemize}[leftmargin=2em]
  \item $L_u = \{a_{u,1}, a_{u,2}, \dots, a_{u,m_u}\}$ is the set of catalog
        anime the user has explicitly marked as \emph{liked}. Each
        $a_{u,j}$ identifies a row in the catalog and carries a
        precomputed embedding $\mathbf{e}_{u,j} \in \mathbb{R}^{384}$
        already stored in Postgres.
  \item $T_u \in \Sigma^{*}$ is the user's optional free-text mood
        description (\textit{e.g.\,``I want something slow-paced and
        melancholic, ideally short''}).
\end{itemize}

At submission time the user supplies at least one of $L_u$ or $T_u$;
empty submissions are rejected.

The catalog is a fixed set $C = \{c_1, \dots, c_N\}$ with $N \approx 38{,}000$.
Every $c \in C$ has a precomputed embedding
$\mathbf{e}_c \in \mathbb{R}^{384}$ produced offline by
\texttt{sentence-transformers/all-MiniLM-L6-v2} over a structured text blob
(title + tags + studios + a short description). All embeddings are
$\ell_2$-normalized, so dot products are cosine similarities in $[-1, 1]$.

\section{Per-user signals}

Once all users have submitted, we derive three quantities per user.

\paragraph{Liked-item matrix.}
We fetch each $\mathbf{e}_{u,j}$ from the catalog (no new model calls)
and stack them:
\[
  \mathbf{L}_u \in \mathbb{R}^{m_u \times 384}.
\]

\paragraph{Text vector.}
If $T_u \neq \varnothing$, we compute
$\mathbf{t}_u = \texttt{embed}(T_u) \in \mathbb{R}^{384}$ on the fly.
This is the only embedding-model call at request time.

\paragraph{Retrieval query.}
A single 384-dim vector that summarizes the user's taste, used only
for candidate retrieval:
\[
  \mathbf{q}_u =
  \begin{cases}
    \mathrm{normalize}\!\left(\tfrac{1}{m_u}\sum_{j} \mathbf{e}_{u,j}\right) & \text{if } L_u \neq \varnothing, \\[4pt]
    \mathbf{t}_u & \text{otherwise (text-only fallback).}
  \end{cases}
\]

For users who supply only liked items, we also reuse $\mathbf{q}_u$ as
their ``text vector'' so the scoring formula stays well-defined without
privileging text-only users.

\section{Pipeline}

\subsection{Step 1: Per-user retrieval}

For each user $u$, we run a pgvector cosine-distance search against the
catalog under the host's hard constraints (year range, allowed media
types) and exclude any anime already in $\bigcup_v L_v$:
\[
  R_u \;=\; \mathrm{topk}_{c \in C \setminus \bigcup_v L_v}\bigl(
    \mathbf{q}_u \cdot \mathbf{e}_c
  \bigr),\qquad |R_u| = 100.
\]
The exclusion ensures we never recommend back to a user something
\emph{any} room member already flagged as liked.

\subsection{Step 2: Group candidate pool}

We union and deduplicate:
\[
  \mathcal{R} \;=\; \bigcup_{u=1}^{n} R_u,
\]
giving a pool of typically 200--500 unique candidates. Stack their
embeddings into $\mathbf{E} \in \mathbb{R}^{|\mathcal{R}| \times 384}$.

\subsection{Step 3: GroupFit pos+text scoring}

For each candidate $i \in \mathcal{R}$ we compute
\begin{equation}
  \mathrm{score}(i) \;=\;
  \underbrace{(1 - \alpha)\,\min_{u}\;\max_{j}\;\bigl(\mathbf{e}_i \cdot \mathbf{e}_{u,j}\bigr)}_{\text{positive term: per-user fairness floor}}
  \;+\;
  \underbrace{\alpha\,\frac{1}{n}\sum_{u}\;\bigl(\mathbf{e}_i \cdot \mathbf{t}_u\bigr)}_{\text{text term: average mood alignment}},
  \label{eq:groupfit}
\end{equation}
with $\alpha = 0.30$ chosen by the offline benchmark (see
\texttt{writeup.tex}).

The two terms answer complementary questions:
\begin{itemize}[leftmargin=2em]
  \item \textbf{Positive term.} For each user $u$, find their
        \emph{best-matching} liked item against candidate $i$
        ($\max_j$). Then take the \emph{worst} of those across users
        ($\min_u$). A high score here means \emph{everyone} has at least
        one liked anime that $i$ resembles --- no user is ignored.
  \item \textbf{Text term.} Average alignment with each user's mood
        description. This is symmetric and forgiving; it nudges the
        ranking but cannot dominate when $\alpha$ is small.
\end{itemize}

In practice we compute the sims in two batched matrix multiplies,
$\mathbf{L}_u \mathbf{E}^\top$ and $\mathbf{T} \mathbf{E}^\top$, then
reduce with \texttt{max}, \texttt{min}, and \texttt{mean}.

\subsection{Step 4: Clustering}

We run a manual K-means implementation
(\texttt{api/app/ml/kmeans.py}) over $\mathbf{E}$, choosing $k$ in
$\{2, \dots, 6\}$ by silhouette score. Each candidate $i$ is assigned a
cluster label $z(i) \in \{1, \dots, k\}$. Clusters typically reflect
thematic groupings in embedding space (\textit{e.g.\,} action-shounen,
slice-of-life, psychological) although they are not literally
genre-labeled.

\subsection{Step 5: Per-cluster ranking and final list}

Within each cluster $\kappa$ we rank by GroupFit score and keep the
top 5 for the UI's clustered shortlist view:
\[
  \mathrm{Cluster}_\kappa = \mathrm{top}_5\bigl(\{i : z(i) = \kappa\},\,
                                                 \mathrm{score}\bigr).
\]
For voting we then take the top 2 per cluster, deduplicate across
clusters, and re-sort globally by score:
\[
  F \;=\; \mathrm{sort}_{\downarrow\mathrm{score}}\Bigl(
    \bigcup_{\kappa=1}^{k} \mathrm{top}_2\bigl(\{i : z(i) = \kappa\},\,
                                                \mathrm{score}\bigr)
  \Bigr).
\]
$F$ is the final voting list shown to all users.

\subsection{Step 6: Voting and result}

Each user submits a subset $V_u \subseteq F$. The room's final result
ranks $F$ first by total vote count
$v(i) = |\{u : i \in V_u\}|$, then by GroupFit score, then by title.
Items tied for the top vote count are flagged as winners.

\section{Cost summary}

The whole compute is dominated by I/O and a few small matrix
multiplications:
\begin{itemize}[leftmargin=2em]
  \item \textbf{Embedding model calls at request time:} only the
        $T_u$ texts (\(\leq n\) short strings), $\sim 50$ ms total
        on warm cache.
  \item \textbf{Catalog embeddings:} all reused from the offline
        preprocessing run; no model calls.
  \item \textbf{Vector search:} $n$ pgvector cosine-distance queries,
        each returning 100 rows under SQL hard-constraint filters.
  \item \textbf{Scoring \& clustering:} $O(|\mathcal{R}| \cdot 384)$
        and $O(|\mathcal{R}| \cdot k \cdot 384)$ respectively;
        milliseconds in practice.
\end{itemize}

The first compute after server start incurs an additional one-time
$\sim 3$--$6$ s cost to load the
\texttt{all-MiniLM-L6-v2} weights into memory; subsequent computes
hit a warm model.

\end{document}